\section{The certificate theorem}
\label{sec:certificate}

Only a family of hyperplanes approaching $E$ is needed to answer
Question~\ref{question:bds}. We first make this weaker assumption precise.

\begin{definition}[Asymptotic hyperplane convexity]\label{def:ahc}
The map $f^h$ has \emph{asymptotic hyperplane convexity} (AHC) if, for
every $\alpha\ne0$ and every $R\in\R$, there exists $s>R$ such that
$f^h(H_{\alpha,s})$ is convex.
\end{definition}

Equivalently, for each $\alpha\ne0$ there is a strictly increasing
sequence $(s_k)$ tending to $+\infty$ with all these images convex.
Indeed, the sequence condition implies the unbounded-level condition;
conversely, choose a good $s_1>1$ and recursively a good
$s_{k+1}>\max\{s_k,k+1\}$. HHC implies AHC because
$H_{\alpha,s}$ is a linear hyperplane for every $s$.
This assumption concerns the images themselves and requires no closedness.

\begin{theorem}[Convex aggregation certificate]\label{thm:certificate}
Suppose $S\ne\varnothing$ and $f^h$ has AHC. Then
\begin{equation}\label{eq:main-certificate}
 \conv S\ne\R^n
 \quad\Longleftrightarrow\quad
 \exists\lambda\geq0:\quad
 A_\lambda\succeq0,\qquad (A_\lambda,b_\lambda)\ne(0,0).
\end{equation}
The certificate-to-proper-hull implication holds without either assumption.
\end{theorem}

\begin{corollary}[Conjecture 3.3 of Blekherman--Dey--Sun]
\label{cor:hhc-certificate}
If $S\ne\varnothing$ and $f^h$ has HHC, then
\eqref{eq:main-certificate} holds.
\end{corollary}

The corollary follows from HHC implying AHC. It answers the cited
conjecture for every finite number of inequalities, including the case
where nonzero convex certificates exist but all their quadratic and linear
coefficients might vanish. The proof below shows that this case forces a
whole-space hull. The theorem also isolates a weaker sufficient hypothesis;
we attach no novelty claim to the separation or compactness facts it uses.

\subsection{Hyperplanes beyond a supporting halfspace}

\begin{lemma}[No common strict negative quadratic direction]
\label{lem:negative-direction}
If $\conv S\ne\R^n$, then there is no $v\in\R^n$ satisfying
$v\transpose A_i v<0$ for every $i$. Equivalently, $E\cap S^h=\varnothing$.
\end{lemma}
\begin{proof}
If such a $v$ existed, then for any fixed $y\in\R^n$,
\[
 f_i(y+rv)=r^2v\transpose A_i v
       +2r(v\transpose A_i y+b_i\transpose v)+f_i(y)
       \longrightarrow-\infty\qquad (|r|\longrightarrow\infty).
\]
There are finitely many constraints. Hence both $y+rv$ and $y-rv$
belong to $S$ for all sufficiently large positive $r$, and their midpoint
$y$ belongs to $\conv S$. This holds for every $y$, a contradiction.
Finally, $f_i^h(v,0)=v\transpose A_i v$ proves the equivalence.
\end{proof}

This lemma concerns a simultaneous strict sign condition on the leading
forms. It is not a characterization of the recession cone of $S$.

\begin{lemma}[Sweeping hyperplanes]\label{lem:sweeping}
Suppose $S\ne\varnothing$ and $\conv S\ne\R^n$. There are
$\alpha\ne0$ and $\beta\in\R$ such that
\begin{equation}\label{eq:strict-support}
 \alpha\transpose x<\beta\qquad(x\in S).
\end{equation}
For any such $\alpha,\beta$, one has
$H_{\alpha,s}\cap S^h=\varnothing$ for every $s\geq\beta$.
\end{lemma}
\begin{proof}
The set $S$ is open, and its convex hull is open as well: in a finite
convex combination, one point with positive weight can vary in an open
ball, giving an open ball around the combination. Separate a point
$y\notin\conv S$ from this nonempty open convex set. Finite-dimensional
separation supplies $\alpha\ne0$ with
$\alpha\transpose x\leq\alpha\transpose y$ throughout $\conv S$.
Equality at a point of that open set would be contradicted by a small
move in direction $\alpha$. Thus \eqref{eq:strict-support} holds with
$\beta=\alpha\transpose y$.

If $(x,t)\in H_{\alpha,s}\cap S^h$ and $t=0$, then
Lemma~\ref{lem:negative-direction} gives a contradiction. If $t\ne0$,
the identity $f_i^h(x,t)=t^2f_i(x/t)$ gives $x/t\in S$, whereas
$\alpha\transpose(x/t)=s\geq\beta$, again a contradiction.
This reasoning applies to either sign of $t$.
\end{proof}

\begin{lemma}[Separation on a hyperplane]\label{lem:hyperplane-certificate}
Suppose $s\ne0$, the set $f^h(H_{\alpha,s})$ is convex, and
$H_{\alpha,s}\cap S^h=\varnothing$. There exists $\lambda\geq0$ with
$\sum_i\lambda_i=1$ such that, for every $x\in\R^n$,
\begin{equation}\label{eq:sweep-inequality}
 0\leq x\transpose A_\lambda x
       +\frac{2}{s}(\alpha\transpose x)(b_\lambda\transpose x)
       +\frac{c_\lambda}{s^2}(\alpha\transpose x)^2.
\end{equation}
\end{lemma}
\begin{proof}
The convex image $C=f^h(H_{\alpha,s})$ contains zero and is disjoint
from the open negative orthant $N=-\R^m_{++}$. Separation of these two
convex sets gives a nonzero $\lambda$ and a real $a$ with
$\lambda\transpose z\leq a\leq\lambda\transpose y$ for $z\in N$,
$y\in C$. If some $\lambda_i<0$, sending $z_i$ to $-\infty$ with
the other coordinates fixed contradicts the upper bound. Thus
$\lambda\geq0$, and $\sup_{z\in N}\lambda\transpose z=0$.
Since $0\in C$, necessarily $a=0$, and
$\lambda\transpose f^h\geq0$ on $H_{\alpha,s}$.
Divide $\lambda$ by its positive coordinate sum. Substituting
$(x,t)=(x,\alpha\transpose x/s)$ gives \eqref{eq:sweep-inequality}.
Closedness of $C$ was not used.
\end{proof}

\subsection{A nonzero limit in coefficient space}

The relevant cone retains the quadratic and linear coefficients together:
\begin{equation}\label{eq:coefficient-cone}
 P=\cone\{(A_i,b_i):1\leq i\leq m\}
   =\{(A_\lambda,b_\lambda):\lambda\geq0\}
   \subseteq\Sym^n\times\R^n.
\end{equation}
We equip this space with
$\|(A,b)\|=(\|A\|_F^2+\|b\|^2)^{1/2}$.
Closedness of this particular cone is essential: a limit in its closure
must correspond to actual aggregation weights. The following elementary
fact supplies that property.

\begin{lemma}[Finite coefficient cones are closed]\label{lem:finite-cone}
The cone generated by a finite family in a finite-dimensional real vector
space is closed.
\end{lemma}
\begin{proof}
Every conic combination can be represented using a linearly independent
subfamily. To see this, choose a representation with the fewest positive
coefficients. If its active generators are dependent, take a nonzero
relation $\sum_i d_i v_i=0$, reversing its sign so some $d_i>0$.
For the active coefficients $a_i>0$, set
$\tau=\min_{d_i>0}a_i/d_i$. Then all $a_i-\tau d_i$ are nonnegative,
at least one vanishes, and the represented point is unchanged, a
contradiction. For a linearly independent subfamily, the coefficient
map is a linear homeomorphism onto its span; that span is closed, so
the image of the nonnegative orthant is closed. There are only finitely
many subfamilies. The original cone is their finite union and is closed.
The empty subfamily accounts for the zero vector.
\end{proof}

\begin{proof}[Proof of Theorem~\ref{thm:certificate}]
First suppose a nontrivial convex certificate $\lambda$ exists.
Its nonconstant coefficient pair implies $\lambda\ne0$, so
$f_\lambda<0$ on $S$. Positive semidefiniteness makes $f_\lambda$
convex, whence $\conv S\subseteq\{f_\lambda<0\}$.
This sublevel set is proper. Indeed, if $A_\lambda\ne0$, it has a
unit eigenvector $v$ with eigenvalue $\mu>0$, and
$f_\lambda(rv)=\mu r^2+2r b_\lambda\transpose v+c_\lambda\to+\infty$
as $r\to+\infty$. If $A_\lambda=0$, then $b_\lambda\ne0$ and
$f_\lambda(r b_\lambda)=2r\|b_\lambda\|^2+c_\lambda\to+\infty$.
This proves the easy implication without using nonemptiness or AHC.

For the converse, suppose $S\ne\varnothing$, $f^h$ has AHC, and
$\conv S\ne\R^n$. Fix $x_0\in S$, and choose $\alpha,\beta$
from Lemma~\ref{lem:sweeping}. AHC gives positive levels
$s_k\to+\infty$ with $s_k\geq\beta$ and convex images
$f^h(H_{\alpha,s_k})$. Lemmas~\ref{lem:sweeping}
and~\ref{lem:hyperplane-certificate} give nonzero normalized
weights $\lambda^k\geq0$ satisfying \eqref{eq:sweep-inequality}.
Write $(A_k,b_k,c_k)=(A_{\lambda^k},b_{\lambda^k},c_{\lambda^k})$.
Strict feasibility gives the crucial upper bound
\begin{equation}\label{eq:constant-elimination}
 c_k<-x_0\transpose A_kx_0-2b_k\transpose x_0
       =:L(A_k,b_k),
\end{equation}
where $L$ is a fixed continuous linear functional on coefficient space.
The multiplier of $c_k$ in \eqref{eq:sweep-inequality} is a square.
Replacing $c_k$ by its upper bound therefore yields
\begin{equation}\label{eq:eliminated-sweep}
 0\leq x\transpose A_kx
      +\frac{2}{s_k}(\alpha\transpose x)(b_k\transpose x)
      +\frac{L(A_k,b_k)}{s_k^2}(\alpha\transpose x)^2
      \qquad(x\in\R^n).
\end{equation}
Moreover, $(A_k,b_k)\ne(0,0)$: otherwise \eqref{eq:constant-elimination}
gives $c_k<0$, and the original inequality
\eqref{eq:sweep-inequality} at $x=\alpha$ reads
$0\leq c_k\|\alpha\|^4/s_k^2<0$.

Set $\rho_k=\|(A_k,b_k)\|>0$ and
$(\widehat A_k,\widehat b_k)=(A_k,b_k)/\rho_k$.
These pairs lie in the unit section of $P$, which is compact by
Lemma~\ref{lem:finite-cone}. Passing to a subsequence, let
\[
 (\widehat A_k,\widehat b_k)\longrightarrow(A_*,b_*)\in P,
 \qquad \|(A_*,b_*)\|=1.
\]
Divide \eqref{eq:eliminated-sweep} by $\rho_k$. By linearity of $L$,
its right-hand side becomes
\[
 x\transpose\widehat A_kx
 +\frac{2}{s_k}(\alpha\transpose x)(\widehat b_k\transpose x)
 +\frac{L(\widehat A_k,\widehat b_k)}{s_k^2}(\alpha\transpose x)^2.
\]
For every fixed $x$ the last two terms tend to zero. Consequently
$x\transpose A_*x\geq0$ for every $x$. This conclusion uses one
coefficient subsequence for all $x$; no point-dependent subsequence is
chosen. The symmetric matrix $A_*$ is positive semidefinite.
Since $(A_*,b_*)\in P$, there are actual weights $\lambda^*\geq0$
with $(A_{\lambda^*},b_{\lambda^*})=(A_*,b_*)$.
Their coefficient pair has norm one, so they give the required
nontrivial convex certificate.
\end{proof}

\begin{remark}[What the limit proves]\label{rem:limit-scope}
Normalizing the weights in the simplex alone can give a negative constant
aggregation in the limit. Normalizing the nonconstant coefficient pair
prevents this loss of information. Bound~\eqref{eq:constant-elimination}
removes the constant before normalization; no bound on $c_k/\rho_k$ is
assumed. The final weights need not be a limit of the original weights,
or of their rescalings. Their existence follows from closedness of the
finitely generated cone $P$.
\end{remark}

\begin{remark}[Convexity along one normal]\label{rem:one-normal}
The hard implication needs convexity of $f^h(H_{\alpha,s})$ at
arbitrarily large levels for only one $\alpha$ satisfying
\eqref{eq:strict-support}. AHC is a convenient condition that can be
stated without knowing a supporting normal. For fixed $\alpha$, the
hyperplanes $H_{\alpha,s}$ approach $E$ as $s\to+\infty$;
convexity of $f^h(E)$ alone is not the hypothesis used in the proof.
There is no restriction beyond $n,m\geq1$. For one constraint,
a negative quadratic direction gives a whole-space hull by
Lemma~\ref{lem:negative-direction}; the remaining cases follow from
the easy implication or a negative constant. For two constraints,
Dines' theorem makes HHC automatic, recovering the known certificate
case discussed in Section~\ref{sec:setting}.
\end{remark}

\subsection{A quantitative form of the obstruction}
\label{sec:quantitative-obstruction}

For completeness, we give a second way to exclude degeneration to
constant certificates. It supplies a uniform bound on the sweep level
under the hypothesis that all convex certificates are trivial. The
cone estimate itself is a standard compactness argument.

\begin{lemma}[Uniform cone separation]\label{lem:uniform-cones}
Let $P,D$ be closed cones containing zero in a finite-dimensional normed
space, with $P\cap D=\{0\}$. Then some $\kappa>0$ satisfies
\[
 \dist(p,D)\geq\kappa\|p\|\qquad(p\in P).
\]
Neither cone is required to be convex.
\end{lemma}
\begin{proof}
If $P=\{0\}$, any positive $\kappa$ works. Otherwise the continuous
function $p\mapsto\dist(p,D)$ has a positive minimum $\kappa$ on the
compact set $P\cap\{\|p\|=1\}$: a zero distance would put a unit
point of $P$ in the closed set $D$. Scaling both cones proves the bound
for $p\ne0$, and the assertion at zero follows from $0\in D$.
\end{proof}

Apply the lemma to $P$ in \eqref{eq:coefficient-cone} and
$D=\Sym^n_+\times\R^n$. The condition $P\cap D=\{0\}$ is exactly
the assertion that every convex certificate is trivial. If $A_-$ is
the positive semidefinite negative part of a symmetric matrix $A$, then
\begin{equation}\label{eq:distance-psd}
 \dist((A,b),D)=\|A_-\|_F
 \leq\sqrt n\max\{0,-\lambda_{\min}(A)\}.
\end{equation}
For the equality, diagonalize $A$ orthogonally. In that basis every
positive semidefinite approximant has nonnegative diagonal entries;
the sum of squared errors is at least the sum of squares of the
negative eigenvalues of $A$. Replacing those eigenvalues by zero
achieves this bound, and the $b$ component can be left unchanged.
The inequality follows by bounding each negative eigenvalue in
absolute value by $\max\{0,-\lambda_{\min}(A)\}$.

\begin{proposition}[Bound on a sweep certificate]\label{prop:sweep-bound}
Suppose every convex certificate is trivial, and let $\kappa$ be as
in Lemma~\ref{lem:uniform-cones} for the preceding cones. If
$\lambda\geq0$, $c_\lambda<0$, $s>0$, $\alpha\ne0$, and
\eqref{eq:sweep-inequality} holds, then
\begin{equation}\label{eq:sweep-level-bound}
 s\leq\frac{2\sqrt n\|\alpha\|}{\kappa}.
\end{equation}
\end{proposition}
\begin{proof}
The inequality at $x=\alpha$ shows that
$\rho=\|(A_\lambda,b_\lambda)\|$ cannot vanish.
The cone estimate and \eqref{eq:distance-psd} give
$\lambda_{\min}(A_\lambda)\leq-\kappa\rho/\sqrt n$.
Choose a unit eigenvector $u$ for this eigenvalue. Evaluating
\eqref{eq:sweep-inequality} at $u$, discarding the nonpositive
constant term, and using $\|b_\lambda\|\leq\rho$ gives
\[
 0\leq-\frac{\kappa}{\sqrt n}\rho
          +\frac{2\|\alpha\|}{s}\rho.
\]
Division by $\rho>0$ proves \eqref{eq:sweep-level-bound}.
\end{proof}

This gives an alternative completion of the hard implication. Under
its hypotheses, suppose every convex certificate is trivial and choose
the normalized sweep weights $\lambda^k$ as before. Compactness of
the simplex gives a subsequence converging to $\bar\lambda$, with
$\bar\lambda\geq0$ and $\sum_i\bar\lambda_i=1$.
The coefficients $b_k,c_k$ are bounded, so letting $k\to\infty$ in
\eqref{eq:sweep-inequality} shows $A_{\bar\lambda}\succeq0$.
Triviality then gives $A_{\bar\lambda}=0$, $b_{\bar\lambda}=0$,
and $c_{\bar\lambda}<0$ by strict feasibility. Thus $c_k<0$
eventually along this subsequence. Proposition~\ref{prop:sweep-bound}
contradicts $s_k\to+\infty$.
The constant $\kappa$ depends only on the original quadratic and linear
coefficients; this argument neither bounds aggregation weights nor
provides a procedure for verifying AHC.
